\chapter*{Abstract}
\addcontentsline{toc}{chapter}{Abstract}
Claims that a new training idea ``saves compute'' are easy to make and hard to trust: baselines are tuned less than the new method, the cost of the extra machinery is not counted, and a small tuning budget gives noisy winners. This report tests two ideas for cheaper neural-network training under one strict protocol. The setting is deliberately small and fully controlled: a four-layer multilayer perceptron on MNIST, cost measured in counted floating-point operations (FLOPs) to reach a fixed validation accuracy, every arm tuned by the same 25-trial random search, paired seeds, and hypotheses, margins and decision rules locked in git before the confirmatory runs.

\textbf{Study~1} asks whether an auxiliary ``meeting loss'', a simple difference-target-propagation target applied at the middle layer next to ordinary backpropagation, reaches the target accuracy with fewer FLOPs. It does not. Against the tuned baseline it needed $36.9\%$ ($97.5\%$ interval $[26.4,47.4]$) more FLOPs to reach $0.97$ and $35.5\%$ ($[26.9,44.4]$) more to reach $0.975$ over 20 paired seeds; a deep-supervision control showed no gain either. Follow-up checks found that the meeting direction is well aligned with the task gradient (cosine $0.75$--$0.80$) and that no tested strength of it sped training up by $10\%$ or more, and two further checks, an oracle target with an exact inverse next to backpropagation and a replacement test with backpropagation cut at the middle layer, found no gain either (all exploratory): the signal does not help whether it is learned or exact, small or strong, and the learned inverse trains the lower layers no better than a frozen random one. This is consistent with, but does not prove, that the signal is redundant with backpropagation in this network.

\textbf{Study~2} turns a measurement from that analysis (a small share of examples carries most of the gradient) into a selective backward pass that back-propagates only the highest-loss $30\%$ of each batch. In a second pre-registered test it cost $35.1\%$ and $38.4\%$ fewer FLOPs than the tuned baseline, and $30.2\%$ and $35.5\%$ fewer than a baseline tuned for a smaller batch, which excludes a pure batch-size explanation. \textbf{Study~3} pre-registers a leaner variant (keep $20\%$, no full-backward warm-up) and finds $48.5\%$ and $48.8\%$ fewer FLOPs than the baseline, $22.2\%$ and $19.1\%$ fewer than the first selective variant, and $46.5\%$ and $53.6\%$ fewer than a tuned reimplementation of Selective-Backprop.

The protocol also produced two methodological results: two tunings of the \emph{same} baseline that tied at selection differed by up to $34\%$ on fresh seeds, and FLOPs and measured time disagree (a per-step time ratio of $0.82$ at one CPU thread against a FLOPs ratio of $0.61$). All results are for one small network on one dataset and count FLOPs, not wall-clock time to target; the report states these limits, lists every post-hoc decision, and proposes the experiments needed to lift them.
\cleardoublepage
